\documentclass{article}
\usepackage{graphicx}
\usepackage{authblk}
\usepackage{booktabs}
\usepackage{float}
\usepackage[brazil]{babel}
\usepackage[utf8]{inputenc}
\usepackage[T1]{fontenc}
\usepackage{indentfirst}
\begin{document}
\begin{titlepage}
	\begin{center}
	
	\vspace{115pt}
    \textbf{\Huge{Projeto Programação Concorrente e Paralela}}\\
        
	\vspace{115pt}
    Carlos Eduardo Nogueira Silva \\
    Felipe Gomes da Silva \\
    Gabriel Martins Brum \\
    Luis Henrique Salomão Lobato \\
	\end{center}
	
	\vspace{1cm}
	\begin{center}
		\vspace{\fill}
    \large{Setembro, 2025} 
	\end{center}
\end{titlepage}
\date{September 2025}




\section{Introdução}

Sistemas  de equações lineares (\textit{Ax=b} podem ser resolvidos de diferentes formas, sendo que um dos possíveis métodos de convergência para identificar soluções é o método de Jacobi, uma abordagem de iterações. 

Neste trabalho foram medidos os tempos de execução a partir da implementação do método de Jacobi tanto sequencial, quanto em paralelo através do uso da ferramenta openMP. 

\section{Metodologia}
\subsection{Cálculo}
O método iterativo de Jacobi foi representado algebricamente através da composição de uma matriz diagonal e as matrizes superiores e inferiores da matriz A. Através desta definição, podê-se chegar em uma solução para um sistema linear através da aplicação da equação \ref{eq:jacobi}. 

O critério de parada adotado foi o de que diferença máxima entre $X_{i}^{k+1}$ e $X_{i}^{k}$ fosse de $10^{-5}$.

\begin{equation}
\label{eq:jacobi}
    x_{i}^{k+1} = \frac{1}{a_{ii}}\left(b_i - \sum_{j \neq i} a_{ij}x_{j}^{k}\right)
\end{equation}

\subsection{Testes}


Foram testadas tanto a implementação sequencial do método, quanto como a execução paralela do algoritmo, com diferentes parâmetros de ajuste. 

O contexto de aplicação dos testes respeitou matrizes A com dimensões de 2000 x 2000 e por consequência, vetores b com 2000 constantes, de modo que os valores utilizados foram de números reais de até 4 casas decimais, respeitando o intervalo de -500 a 500. 

Os parâmetros avaliados ao longo dos testes com a versão openMP foram o número de threads de execução variando de 2 a 8, o tipo de escalonador de threads (\textit{static}, \textit{guided} e \textit{dynamic})  e o número de \textit{chunks} (0 ou 25).  

Os testes foram executados em diferentes máquinas com diferentes processadores, visando evitar qualquer tipo de viés na análise de desempenho dos métodos e aumentar o contexto de exploração proposto. 


\section{Resultados}
Na Tabela \ref{tab:processor_specs} estão detalhados os processadores utilizados para a executar dos programas sequencial e openMP, seguido pelo número de núcleos e \textit{threads} dos mesmos.


\begin{table}[h]
    \centering
    \caption{Especificações dos processadores onde os algoritmos foram testados.}
    \label{tab:processor_specs}
    \begin{tabular}{ccc}
    \toprule
    \textbf{Processador} & \textbf{Núcleos} & \textbf{Threads} \\
    \midrule
    Intel Core i3-10100F & 4 & 8 \\
    Intel Core i7-10750H & 6 & 12 \\
    AMD Ryzen 7 5700u & 8 & 16 \\
    AMD Ryzen 5 3500u & 4 & 8 \\
    \bottomrule
    \end{tabular}
\end{table}
Assim, encontram-se nas Tabelas  \ref{tab:results}, \ref{tab:results2}, \ref{tab:results3} e \ref{tab:results4} os resultados das medições do tempo de execução dos algoritmos em cada processador. 

\begin{table}[H]
    \caption{Resultados da execução dos testes em um i3-10100F.}
    \centering
    \label{tab:results}
    \begin{tabular}{cccc}
    \toprule
    \textbf{Threads} & \textbf{Schedule} & \textbf{Chunk} & \textbf{Tempo de Execução} \\
    \midrule
    1 (sequencial) & N/A & N/A & 0.1024s\\
    \addlinespace
    2 & Estático & N/A & 1.1503s \\
    4 & Estático & N/A & 2.6165s \\
    8 & Estático & N/A & 3.5367s \\
    \addlinespace
    2 & Estático & 25 &  1.1271s\\
    4 & Estático & 25 & 2.3659s\\
    8 & Estático & 25 &  2.6648s\\
    \addlinespace
    2 & Dinâmico & 25 &  1.1453s\\
    4 & Dinâmico & 25 & 2.1047s\\
    8 & Dinâmico & 25 &  2.8254s\\
    \addlinespace
    2 & Guided & 25 &  1.0959s\\
    4 & Guided & 25 &  2.6802s\\
    8 & Guided & 25 &  3.1195s\\
    
    \bottomrule
    \end{tabular}
\end{table}

\begin{table}[H]
    \caption{Resultados da execução dos testes em um i7-10750H.}
    \centering
    \label{tab:results3}
    \begin{tabular}{cccc}
    \toprule
    \textbf{Threads} & \textbf{Schedule} & \textbf{Chunk} & \textbf{Tempo de Execução} \\
    \midrule
    1 (sequencial) & N/A & N/A & 0.4097s \\
    \addlinespace
    2 & static & 1 & 4.3661s \\
    4 & static & 1 & 8.7037s \\
    8 & static & 1 & 11.7228s \\
    \addlinespace
    2 & static & 25 & 2.9943s \\
    4 & static & 25 & 7.0912s \\
    8 & static & 25 & 10.9250s \\
    \addlinespace
    2 & dynamic & 25 & 4.0930s \\
    4 & dynamic & 25 & 8.7469s \\
    8 & dynamic & 25 & 11.4043s \\
    \addlinespace
    2 & guided & 25 & 4.0025s \\
    4 & guided & 25 & 9.1439s \\
    8 & guided & 25 & 13.2075s \\
    \bottomrule
    \end{tabular}
\end{table}


\begin{table}[H]
    \caption{Resultados da execução dos testes em um AMD Ryzen 7 5700u.}
    \centering
    \label{tab:results2}
    \begin{tabular}{cccc}
    \toprule
    \textbf{Threads} & \textbf{Schedule} & \textbf{Chunk} & \textbf{Tempo de Execução} \\
    \midrule
    1 (sequencial) & N/A & N/A & 0.1277s\\
    \addlinespace
    2 & Estático & 1 & 1.0148s \\
    4 & Estático & 1 & 2.0533s \\
    8 & Estático & 1 & 2.6728s \\
    \addlinespace
    2 & Estático & 25 & 1.0125s\\
    4 & Estático & 25 & 2.2540s\\
    8 & Estático & 25 & 2.6751s\\
    \addlinespace
    2 & Dinâmico & 25 & 1.0092s\\
    4 & Dinâmico & 25 & 2.0855s\\
    8 & Dinâmico & 25 & 2.6796s\\
    \addlinespace
    2 & Guided & 25 & 1.0120s\\
    4 & Guided & 25 & 2.1032s\\
    8 & Guided & 25 & 2.6918s\\
    
    \bottomrule
    \end{tabular}
\end{table}


\begin{table}[H]
\caption{Resultados da execução dos testes em um AMD Ryzen 5 3500u.}
 \centering
 \label{tab:results4}
 \begin{tabular}{cccc}
 \toprule
 \textbf{Threads} & \textbf{Schedule} & \textbf{Chunk} & \textbf{Tempo de Execução} \\
 \midrule
 1 (sequencial) & N/A & N/A & 0.1515s\\
 \addlinespace
 2 & static & 1 & 1.0952s \\
 4 & static & 1 & 2.7601s \\
 8 & static & 1 & 4.3522s \\
 \addlinespace
 2 & static & 25 & 1.0758s\\
 4 & static & 25 & 2.7876s\\
 8 & static & 25 & 4.4914s\\
\addlinespace
 2 & dynamic & 25 & 1.0862s\\
4 & dynamic & 25 & 2.8168s\\
 8 & dynamic & 25 & 4.2930s\\
 \addlinespace
 2 & guided & 25 & 1.4712s\\
 4 & guided & 25 & 2.5914s\\
 8 & guided & 25 & 3.8268s\\
 \bottomrule
 \end{tabular}
\end{table}


\section{Conclusão}

A partir dos dados coletados, foi possível concluir que para este caso específico de testes, onde o método de Jacobi acabou por convergir em poucas iterações (neste caso, 7), a implementação sequencial obteve vantagem de forma majoritária, pois o programa openMP acaba por consumir mais tempo para a criação, alocação e escalonamento de threads do que executando em si o cálculo do método. 
\end{document}
